\section{The numerical constraint}
\label{sec:arithmetic}

The geometric argument has produced the bilinear equations
\eqref{eq:N-products}. We now prove that they allow only two values of
$u$. This section is a finite calculation over $\mathbb Q$: its input
is the explicitly defined matrix $A$ and its selected coordinates
$C_1$, and its output is a quadratic equation for $\eta_1$.

\begin{lemma}\label{lem:finite-ratio}
Let $A$, $C_1$ and $N$ be the rational matrices and index set defined
in Section~\ref{sec:characteristic-classes}. Suppose that
$\xi\in\mathbb Q^{11}$ and $\eta\in\mathbb Q^{18}$ satisfy
\[
 \xi_0=\eta_0=1,
 \qquad N(\xi_i\eta_j)_{(i,(j))\in C_1}=0.
\]
Then
\begin{equation}\label{eq:ratio-quadratic}
 32-20\eta_1+3\eta_1^2=0.
\end{equation}
\end{lemma}

We give the finite certificate, including the distinction between
the modular rank bound and the rational identities needed to finish
the proof. The exact recurrences and elimination procedure are
implemented in Appendix~\ref{app:certificate}.

\subsection{Multiplying the constraints}

Write $y_{i,j}=\xi_i\eta_j$ for the selected bilinear coordinates.
Multiplying each equation of $Ny=0$ by $\eta_a\eta_b$, with
$a,b\in J$, gives linear equations in the cubic products
\[
 v_{i,\operatorname{sort}(j,a,b)}=\xi_i\eta_j\eta_a\eta_b,
 \qquad(i,(j))\in C_1.
\]
Here $\operatorname{sort}$ arranges the three indices in nondecreasing
order, identifying equal products.
We will impose those linear equations without requiring an arbitrary
solution $v$ to have this product form. In the resulting larger
solution space, we prove that
\begin{equation}\label{eq:target-functional}
 32v_{0,(0,0,0)}-20v_{0,(0,0,1)}+3v_{0,(0,1,1)}=0.
\end{equation}
For the product vector, $\xi_0=\eta_0=1$ makes the left side
$32-20\eta_1+3\eta_1^2$. Thus a linear statement on cubic
coordinates will prove the desired quadratic relation. The following
spaces make this relaxation precise; the intermediate quadratic
space also makes its computation smaller.

For $n=2,3$, form $C_n$ from $C_{n-1}$ by appending an index $j\in J$
to its tuple and sorting. Repeated coordinates are identified.
Equivalently, $(i,z)\in C_n$ means that $z$ is a size-$n$ multiset
on $J$, and either $h\mid M_i$ or $z$ contains an index of weight
three. Let $\mathcal K_1=\ker N$. Recursively define
$\mathcal K_n\subseteq\mathbb Q^{C_n}$ by requiring, for every
$j\in J$, that the slice
\[
 (v_{i,\operatorname{sort}(z,j)})_{(i,z)\in C_{n-1}}
 \quad\text{belong to }\mathcal K_{n-1}.
\]
Thus $\mathcal K_3$ is the kernel of a single rational matrix:
for each unordered pair $a,b\in J$, allowing $a=b$, impose $N$ on the slice
\[
 \bigl(v_{i,\operatorname{sort}(j,a,b)}\bigr)_{(i,(j))\in C_1}.
\]
There are $\binom{16}{2}=120$ such pairs. In particular,
\begin{equation}\label{eq:geometric-cubic-column}
 v_{i,z}=\xi_i\prod_{j\in z}\eta_j
\end{equation}
belongs to $\mathcal K_3$: every double slice is a scalar multiple
of the given vector in $\ker N$.

Seven monomials $M_i$ contain $h$, and six of the fifteen indices
in $J$ have weight three. Therefore
\[
 |C_n|=7\binom{14+n}{n}
       +4\left(\binom{14+n}{n}-\binom{8+n}{n}\right),
\]
which gives $|C_1|=129$, $|C_2|=1140$ and $|C_3|=6820$.

\subsection{An upper bound from finite-field elimination}

The coefficient recurrences in Section~\ref{sec:characteristic-classes}
give $1785$ nonzero rows and $581$ columns for $A$. Exact elimination
gives a matrix $K$ of size $581\times65$ with $AK=0$.
The selected-row matrix $P$ has size $129\times65$. Choose $N$ in
reduced row echelon form. The following checks give, in particular,
complete bases rather than merely lists of relations:
\begin{center}
\begin{tabular}{c|c|c}
matrix & size & rank modulo $10007$\\\hline
$A$ & $1785\times581$ & $516$\\
$K$ & $581\times65$ & $65$\\
$P$ & $129\times65$ & $60$\\
$N$ & $69\times129$ & $69$
\end{tabular}
\end{center}
The identities $AK=0$ and $NP=0$ hold over $\mathbb Q$.
Every denominator in these matrices is prime to $10007$.
Consequently the first two ranks prove that $K$ is a basis of
$\ker A$, and the last two prove that $N$ is the full left
annihilator of $P$.
Write $\mathcal K_n(\mathbb F_{10007})$ for the spaces defined
by the same slice equations using the reduction of $N$ modulo
$10007$.

Here is the elimination procedure used for the larger spaces.
Suppose that the columns of $Z$ are a basis for the solutions on a union
$S$ of slices, and the columns of $B$ are a basis for the solutions on the
next slice $T$. On the overlap, the two parameter vectors $a,b$
must satisfy
\[
 Z_{S\cap T}a-B_{S\cap T}b=0.
\]
If the columns of $\binom{V_1}{V_2}$ form the kernel of this matrix,
then the columns of
\[
 \begin{pmatrix} ZV_1\\ B_{T\setminus S}V_2\end{pmatrix}
\]
parametrize all solutions on $S\cup T$. This parametrization is
injective: zero on the union is zero on each piece, and the old
parametrizations were injective. Appending a fixed index to a
multiset is itself injective, so there are no hidden identifications
within an individual slice. Ordinary Gaussian elimination and
successive application of this construction give
\begin{equation}\label{eq:finite-field-dimensions}
 \dim\mathcal K_1(\mathbb F_{10007})=60,
 \qquad\dim\mathcal K_2(\mathbb F_{10007})=163,
 \qquad\dim\mathcal K_3(\mathbb F_{10007})=6.
\end{equation}

For the rational conclusion, use the single double-slice matrix
defining $\mathcal K_3$. It has $69\cdot120=8280$ rows and $6820$
columns, and its denominators are units modulo $10007$. Its modular
rank is $6814$, so its rational rank is at least $6814$. Thus
\begin{equation}\label{eq:rational-dimension-upper}
 \dim_{\mathbb Q}\mathcal K_3\le6.
\end{equation}
This argument does not assume that intermediate rational kernel
bases specialize to bases.

\subsection{Six rational columns}

It remains to exhibit six independent rational solutions and show that
the functional giving the desired quadratic vanishes on their span.
These columns are witnesses for the linear system: no assertion that
they arise from varieties is needed. Three have product form, and
three are specified by their small supports. The relaxed system
allows both kinds without imposing the original normalization or
requiring a product representation.
Define the following rational vectors:
\begin{align*}
 B(v)={}&\bigl(1,7-2v/5,45-4v,275-30v,11,65-2v,\\
             &\hspace{29mm}85,17,95-2v,25,35\bigr),\\
 B'={}&(1,11/5,5,11,-1,-1,10,-3,-1,0,14),\\
 H^a={}&\bigl(1,4,16,64,6,24,4,3,12,\\
       &\hspace{29mm}48,192,768,17,68,272,96,9,36\bigr),\\
 H^b={}&\tfrac13\bigl(3,8,21,54,9,24,6,14,37,\\
       &\hspace{29mm}97,252,648,37,97,252,96,16,42\bigr).
\end{align*}
For $(B,H)=(B(0),H^a),(B(1),H^a),(B',H^b)$, let
\[
 v(B,H)_{i,z}=B_i\prod_{j\in z}H_j.
\]
Direct rational substitution gives
$N(B_iH_j)_{C_1}=0$, so these three columns lie in $\mathcal K_3$.
Three further columns $s_1,s_2,s_3$ have the following entries and
are zero elsewhere:
\begin{center}
\begin{tabular}{c|c|rrr}
$i$ & $z$ & $s_1$ & $s_2$ & $s_3$\\\hline
$5$ & $(15,15,15)$ & $3$ & $0$ & $0$\\
$7$ & $(15,15,15)$ & $1$ & $0$ & $0$\\
$9$ & $(15,15,15)$ & $-5$ & $0$ & $0$\\
$7$ & $(10,10,15)$ & $0$ & $2$ & $0$\\
$7$ & $(10,10,10)$ & $0$ & $4$ & $1$\\
$9$ & $(10,10,10)$ & $0$ & $25$ & $0$
\end{tabular}
\end{center}
Applying $N$ to each of their $120$ double slices gives zero over
$\mathbb Q$. These are exact rational equalities, separate from
the finite-field calculation.

The six columns are independent. Indeed, restrict them to the rows
\begin{gather*}
 (0,(0,0,0)),\quad (0,(0,0,1)),\quad (1,(0,0,0)),\\
 (5,(15,15,15)),\quad (7,(10,10,15)),\quad (7,(10,10,10)).
\end{gather*}
The upper-right $3\times3$ block is zero. The upper-left block has
determinant $-8/15$, and the lower-right block is
$\left(\begin{smallmatrix}3&0&0\\0&2&0\\0&4&1\end{smallmatrix}\right)$,
with determinant $6$. The full minor is $-16/5$. Together with
\eqref{eq:rational-dimension-upper}, this proves that the six
displayed columns form a rational basis of $\mathcal K_3$.

\begin{proof}[Proof of Lemma~\ref{lem:finite-ratio}]
The functional in~\eqref{eq:target-functional} vanishes on the six
basis columns. On a displayed product column its value
is $32-20H_1+3H_1^2$, with $H_1=4$ or $8/3$; the sparse columns
are zero at all three coordinates. Apply it to
\eqref{eq:geometric-cubic-column} and use $\xi_0=\eta_0=1$.
\end{proof}

Returning to the evaluation fibre, \eqref{eq:N-products} verifies
the hypothesis of the lemma and $\eta_1=u$. Hence
\[
 (u-4)(3u-8)=0.
\]
It remains to identify which value is compatible with the tangent
map and to recognize the resulting variety of minimal rational
tangents.
